\section{Conclusion}
\label{sec:concl}

We derived backpropagation for a one-hidden-layer perceptron with a sigmoid output and binary cross-entropy, implemented it from scratch in NumPy, and verified it: the analytic gradients agree with centered finite differences to within $2.1\times10^{-11}$. The learning rate, hidden width and regularization strength were chosen from a 27-cell grid by the mean validation loss over ten initialization seeds. The selected configuration ($\eta=0.01$, $n_1=16$, $\lambda=0.01$; mean validation BCE $0.076259$) lies on the edge of the grid on all three axes. Without regularization the validation loss grew with the learning rate, from $0.0796$--$0.0875$ at $\eta=0.01$ to $0.2906$--$0.3179$ at $\eta=0.5$, which is consistent with overfitting; with $\lambda=0.01$ the corresponding ranges were $0.0763$--$0.0826$ and $0.0824$--$0.0898$. The final model stopped at the 5000-epoch cap with a final gradient norm of $0.015658$. On the 86 test examples it reached an accuracy of $0.9884$, a precision of $1.0000$ and a recall of $0.9688$, with a test log loss of $0.0497$; it missed one of the 32 malignant cases.

The study is relevant as an illustration of how a classifier for fine-needle-aspirate features can be built, validated and judged, for example as a screening-support aid. It is an educational model; it is not a diagnostic tool, and the results support no clinical use. As for the method, the MLP matches logistic regression on this split: both make the same single error, and the majority-class rule reaches an accuracy of $0.6279$. The test log losses are $0.0497$ and $0.0610$, two values from 86 examples and one split, and we draw no stronger conclusion from them.

Several extensions were not done and are proposed as future work: (i) external validation on data from another source; (ii) patient-level or cross-validated splits, which would also give an uncertainty estimate for the test metrics; (iii) choosing the decision threshold on validation data to favor recall, since the missed case had $\hat p=0.4619$, just below the threshold $0.5$, a value that is a test-set observation and not a basis for changing the threshold; (iv) a grid that extends beyond the current edge in learning rate, width and regularization strength; and (v) mini-batch training in place of full-batch gradient descent.
